\section{Exact modular transport of the order derivatives}
\label{sec:modular}

Question 7 of \cite{RepoVertical} asks for an exact
\(\tau\leftrightarrow4\pi^2/\tau\) transport, including the
nonmodular terms. At horizontal shift \(a=1\), we give one
holomorphic identity near every nonpositive even order and another
near the Stieltjes order \(s=1\). Their derivatives are specified
by finite polynomial recurrences and complete residue formulas.

This section uses classical Mellin-transform technology. The eta
and Ramanujan specializations, the zeta functional equation, and
the existence of modular relations for generalized Stieltjes
constants are established literature, not novelty claims.
In particular, the all-index modular relations of Dixit,
Sathyanarayana and Sharan \cite{DixitStieltjes} preclude any claim
that modular Stieltjes identities themselves originate here.
The contribution to the present continuation is the explicit
order-dependent kernel, its residue compiler, and its compatibility
with the fixed-cutoff law of Section~\ref{sec:lattice}.

\subsection{One contour identity before coefficient extraction}

Put
\begin{equation}\label{eq:mod-lambda}
 \Lambda(s,\tau)=\sum_{n=1}^{\infty}
       \frac{n^{s-1}}{e^{n\tau}-1},\qquad
 \tau^*=\frac{4\pi^2}{\tau},\qquad
 \alpha=\frac{2\pi}{\tau},\qquad L=\log\alpha.
\end{equation}
The function \(\Lambda\) is entire in \(s\). It is exactly
\({\cal L}_{1-s}(e^{-\tau},e^{-\tau})\), so subtracting the
finite polylogarithm terms in \eqref{eq:root-bracket} recovers
the vertically sampled family.

For \(M\geq0\), let
\begin{align}
 {\cal Q}_M(s,\tau)
  ={}&\frac{\zeta(2-s)}{\tau}
   +\Gamma(s)\zeta(s)\tau^{-s}-\frac12\zeta(1-s)\notag\\
 &+\sum_{j=1}^{M+1}\frac{B_{2j}}{(2j)!}
          \zeta(2-2j-s)\tau^{2j-1}.
 \label{eq:mod-residue-even}
\end{align}
Near \(s=1\), instead put
\begin{equation}\label{eq:mod-residue-one}
 {\cal Q}_*(s,\tau)=
   \frac{\zeta(2-s)}{\tau}
   +\Gamma(s)\zeta(s)\tau^{-s}-\frac12\zeta(1-s).
\end{equation}
The apparently singular terms are combined before their removable
values are taken.

\begin{theorem}[Local modular transport]\label{thm:modular}
Fix \(\tau>0\). If \(M\geq0\) and \(|s+2M|<1/4\), then
\begin{equation}\label{eq:modular-master}
 \boxed{\displaystyle
 \Lambda(s,\tau)={\cal Q}_M(s,\tau)
  +\frac{\alpha^s}{2\pi\ii}
   \int_{\Ree w=3/2}
       \Gamma(w)\zeta(w)\zeta(w+1-s)
       \frac{\cos(\pi w/2)}{\cos(\pi(s-w)/2)}
       (\tau^*)^{-w}\,\dd w.}
\end{equation}
If \(|s-1|<1/4\), the same identity holds with
\({\cal Q}_M\) replaced by \({\cal Q}_*\).
The contour is oriented upward. Both identities are holomorphic
through their central orders, and the integral admits every fixed
order derivative under its sign.
\end{theorem}

\begin{proof}
For \(c>\max(1,\Ree s)\), Mellin inversion and absolute convergence
give
\begin{equation}\label{eq:lambert-mellin}
 \Lambda(s,\tau)=\frac1{2\pi\ii}
       \int_{\Ree z=c}\Gamma(z)\zeta(z)\zeta(z+1-s)
                                  \tau^{-z}\,\dd z.
\end{equation}
In the stated neighborhoods one may take \(c=3/2\) and move the
contour to \(\Ree z=\Ree s-3/2\).
Near \(s=-2M\), the poles crossed are \(z=1,s,0\) and
\(z=1-2j\), \(1\leq j\leq M+1\). Their residues are exactly
\eqref{eq:mod-residue-even}. The negative even poles of Gamma
are canceled by the trivial zeros of \(\zeta(z)\).
Near \(s=1\), only \(z=1,s,0\) are crossed, giving
\eqref{eq:mod-residue-one}.

The two zeta functional equations \cite[Section 25.4]{DLMF} imply
the meromorphic identity
\begin{equation}\label{eq:mod-functional-factor}
 \begin{split}
 \Gamma(z)\zeta(z)\zeta(z+1-s)
  ={}&(2\pi)^{2z-s}\Gamma(s-z)\zeta(s-z)\zeta(1-z)\\
 &\times
 \frac{\cos(\pi(s-z)/2)}{\cos(\pi z/2)}.
 \end{split}
\end{equation}
Substitute \(w=s-z\) on the new contour; the resulting integral
is the one in \eqref{eq:modular-master}.
Stirling's formula gives exponential decay, while the zeta
factors have at most polynomial growth on the fixed vertical
strips. The trigonometric quotient and all its order derivatives
are bounded there on compact parameter sets: the denominator has
no zero on the contour in either stated neighborhood. These facts
justify the contour motion and differentiated integrals.
Away from the central collisions the residue computation is
ordinary. The formulas below explicitly remove those collisions,
or equivalently one may use the identity theorem and the
holomorphic left-hand side.
\end{proof}

At \(s=-2M\), the trigonometric quotient is the constant
\((-1)^M\). Thus the zeroth derivative has the classical
Ramanujan form
\begin{equation}\label{eq:ramanujan-recovered}
 \Lambda(-2M,\tau)=
 {\cal Q}_M(-2M,\tau)+(-1)^M\alpha^{-2M}
                         \Lambda(-2M,\tau^*).
\end{equation}
For \(M=0\), it becomes the eta transformation
\begin{equation}\label{eq:eta-recovered}
 \Lambda(0,\tau)-\Lambda(0,\tau^*)
   =\frac{\zeta(2)}{\tau}
     +\frac12\log\frac{\tau}{2\pi}-\frac{\tau}{24}.
\end{equation}
For example, at \(M=1\),
\begin{align}
 \Lambda(-2,\tau)
 ={}&\frac{\zeta(4)}{\tau}
   -\frac{\zeta(3)}2\left(1+\frac{\tau^2}{4\pi^2}\right)
   +\frac{\zeta(2)\tau}{12}+\frac{\tau^3}{1440}\notag\\
 &-\frac{\tau^2}{4\pi^2}\Lambda(-2,\tau^*).
 \label{eq:zeta3-recovered}
\end{align}
These are useful normalization checks. The positive
\(\tau^3/1440\) sign in the last polynomial is forced by
\(B_4\zeta(0)>0\).

\subsection{A finite polynomial compiler for every derivative}

Define polynomials \(p_j(X)\) by
\begin{equation}\label{eq:mod-kernel-recurrence}
 p_0(X)=1,\qquad
 p_{j+1}(X)=-(1+X^2)p_j'(X)-Xp_j(X).
\end{equation}
They satisfy
\begin{equation}\label{eq:mod-kernel-generator}
 \frac1{\cos v+X\sin v}
      =\sum_{j=0}^{\infty}p_j(X)\frac{v^j}{j!}.
\end{equation}
Indeed, differentiating a shifted secant proves the recurrence.
The first values are
\[
 p_1=-X,\qquad p_2=1+2X^2,\qquad
 p_3=-5X-6X^3,\qquad p_4=5+28X^2+24X^4.
\]

For \(M,j,\ell\geq0\), set
\begin{equation}\label{eq:Tmod}
 {\cal T}_{j,\ell}^{(M)}(x)=
 \frac{(-1)^\ell}{2\pi\ii}
 \int_{\Ree w=3/2}
  \Gamma(w)\zeta(w)\zeta^{(\ell)}(w+1+2M)
  p_j\!\left(\tan\frac{\pi w}{2}\right)x^{-w}\,\dd w.
\end{equation}
Every integral is absolutely convergent. For \(j=0\),
\[
 {\cal T}_{0,\ell}^{(M)}(x)
       =\left.\partial_s^\ell\Lambda(s,x)\right|_{s=-2M}.
\]

\begin{corollary}[All negative-even modular derivatives]\label{cor:mod-jets}
For every \(r\geq0\),
\begin{equation}\label{eq:mod-jets}
 \boxed{\displaystyle
 \partial_s^r\Lambda(-2M,\tau)
 ={\cal Q}_M^{(r)}(-2M,\tau)
  +(-1)^M\alpha^{-2M}
  \sum_{\substack{j,\ell,k\geq0\\j+\ell+k=r}}
       \frac{r!}{j!\ell!k!}
       \left(\frac{\pi}{2}\right)^j
       L^k\,{\cal T}_{j,\ell}^{(M)}(\tau^*).}
\end{equation}
\end{corollary}

\begin{proof}
Put \(s=-2M+u\). The quotient in
\eqref{eq:modular-master} becomes
\[
 \frac{(-1)^M}
 {\cos(\pi u/2)+\tan(\pi w/2)\sin(\pi u/2)}.
\]
Expand this by \eqref{eq:mod-kernel-generator},
\(\alpha^u=e^{Lu}\), and
\(\zeta(w+1+2M-u)\). Normal convergence allows coefficient
extraction. The multinomial coefficient is the ordinary
Leibniz coefficient.
\end{proof}

The terms \(j=0\) are the familiar transformed Lambert
derivatives. The finitely many \(j\geq1\) terms are the exact
nonmodular correction. Omitting them is already wrong at the
first derivative. This is an identity with explicitly specified
ordinary contour integrals, not an assertion that those integrals
always reduce to a finite list of polylogarithm values.

The corresponding critical-order compiler is equally explicit.
Define
\begin{equation}\label{eq:Umod}
 {\cal U}_{j,\ell}(x)=
 \frac{(-1)^\ell}{2\pi\ii}
 \int_{\Ree w=3/2}
    \Gamma(w)\zeta(w)\zeta^{(\ell)}(w)
    Xp_j(-X)x^{-w}\,\dd w,\qquad
    X=\cot\frac{\pi w}{2}.
\end{equation}
Since the trigonometric quotient at \(s=1+u\) is
\[
 \frac{X}{\cos(\pi u/2)-X\sin(\pi u/2)},
\]
we obtain
\begin{equation}\label{eq:critical-mod-jets}
 \partial_s^r\Lambda(1,\tau)
  ={\cal Q}_*^{(r)}(1,\tau)
   +\alpha\!
    \sum_{j+\ell+k=r}\frac{r!}{j!\ell!k!}
          \left(\frac{\pi}{2}\right)^j L^k
                                  {\cal U}_{j,\ell}(\tau^*).
\end{equation}
Thus every critical-order derivative is transported by finitely
many integrals involving \(\zeta(w)\zeta^{(\ell)}(w)\).

\subsection{The colliding residues and their Stieltjes coefficients}

The finite residue terms in the preceding formulas require care.
At zero, let
\begin{equation}\label{eq:mod-e}
 e_j=\left.\frac{\dd^j}{\dd u^j}
                 \left(e^{Lu}\sec\frac{\pi u}{2}\right)\right|_{u=0}.
\end{equation}
Then \(e_0=1\), \(e_1=L\),
\(e_2=L^2+\pi^2/4\), and
\(e_3=L^3+3\pi^2L/4\). The zeta functional equation gives
\[
 \Gamma(s)\zeta(s)\tau^{-s}-\frac12\zeta(1-s)
 =\frac12\zeta(1-s)
           \left(\alpha^s\sec\frac{\pi s}{2}-1\right).
\]
The factor in parentheses vanishes at zero. Multiplying the
full Laurent series, including its pole, yields
\begin{equation}\label{eq:Qzero-jets}
 \begin{split}
 {\cal Q}_0^{(r)}(0,\tau)
 ={}&(-1)^r\left(\frac{\zeta^{(r)}(2)}{\tau}
                  +\frac{\tau}{12}\zeta^{(r)}(0)\right)\\
 &+\frac12\left\{-\frac{e_{r+1}}{r+1}
        +\sum_{j=0}^{r-1}\binom rj\gamma_j e_{r-j}\right\}.
 \end{split}
\end{equation}
The sum is empty at \(r=0\). In particular, the residue at
derivative order \(r\) needs one extra coefficient \(e_{r+1}\).
Replacing the Laurent germ by its regular part would delete this
term.

For \(M\ge1\) there is no Laurent collision at \(s=-2M\).
The same \(e_j\), with \(L=\log\alpha\), give every residue
derivative without separately differentiating a singular Gamma
factor:
\begin{equation}\label{eq:Qnegative-jets}
 \begin{split}
 {\cal Q}_M^{(r)}(-2M,\tau)
 ={}&(-1)^r\left\{\frac{\zeta^{(r)}(2M+2)}{\tau}
    +\sum_{j=1}^{M+1}\frac{B_{2j}\tau^{2j-1}}{(2j)!}
                    \zeta^{(r)}(2M+2-2j)\right\}\\
 &+\frac12\left\{(-1)^M\alpha^{-2M}
      \sum_{\ell=0}^{r}\binom r\ell(-1)^\ell
                \zeta^{(\ell)}(2M+1)e_{r-\ell}
              -(-1)^r\zeta^{(r)}(2M+1)\right\}.
 \end{split}
\end{equation}
Indeed the functional-equation expression immediately preceding
\eqref{eq:Qzero-jets} is a product of two holomorphic functions
at these orders. Ordinary Leibniz differentiation proves
\eqref{eq:Qnegative-jets}.

At \(s=1\), let
\begin{equation}\label{eq:mod-A}
 a_j=\left.\frac{\dd^j}{\dd u^j}
            \left(\Gamma(1+u)\tau^{-u}\right)\right|_{u=0}.
\end{equation}
These are Bell polynomials generated by
\[
 \Gamma(1+u)\tau^{-u}
 =\exp\left(-(\gamma+\log\tau)u+
          \sum_{m=2}^{\infty}\frac{(-1)^m\zeta(m)}m u^m\right).
\]
The two colliding terms at \(s=1+u\) combine as
\[
 \frac1{\tau}
 \left\{\frac{A(u)-1}{u}
       +\sum_{j=0}^{\infty}\frac{\gamma_j u^j}{j!}
                              \bigl(1+(-1)^j A(u)\bigr)\right\},
 \qquad A(u)=\Gamma(1+u)\tau^{-u}.
\]
Consequently,
\begin{equation}\label{eq:Qone-jets}
 \boxed{\displaystyle
 {\cal Q}_*^{(r)}(1,\tau)=
 \frac1{\tau}\left\{\frac{a_{r+1}}{r+1}+\gamma_r+
     \sum_{j=0}^{r}\binom rj(-1)^j\gamma_j a_{r-j}\right\}
       -\frac{(-1)^r}{2}\zeta^{(r)}(0).}
\end{equation}
This is a finite formula in ordinary Stieltjes constants,
zeta values, and logarithms for every \(r\).
The first two cases are
\begin{align}
 {\cal Q}_*(1,\tau)&=\frac{\gamma-\log\tau}{\tau}+\frac14,
 \label{eq:Qone-first}\\
 {\cal Q}_*'(1,\tau)&=
       \frac{\log^2\tau-\gamma^2+\zeta(2)}{2\tau}
                    +\frac12\zeta'(0).
 \label{eq:Qone-second}
\end{align}
Cancellation of \(\gamma_1\) in the latter identity is exact.

\subsection{A self-dual integral evaluation}

At \(\tau=2\pi\), \(M=0\), and \(r=1\), the Lambert
derivatives on the two sides of \eqref{eq:mod-jets} cancel.
Using \eqref{eq:Qzero-jets} gives the explicit identity
\begin{equation}\label{eq:self-dual-integral}
 \boxed{\displaystyle
 \frac1{2\pi\ii}\int_{\Ree w=3/2}
   \Gamma(w)\zeta(w)\zeta(w+1)
      \tan\frac{\pi w}{2}\,(2\pi)^{-w}\,\dd w
 =-\frac{\zeta'(2)}{\pi^2}-\frac{\pi}{8}
                              +\frac{\log(2\pi)}6.}
\end{equation}
The integral is absolutely convergent; no principal value occurs.
Higher self-dual derivatives give finite linear identities among
the integrals \({\cal T}_{j,\ell}^{(0)}(2\pi)\), with complete
Stieltjes coefficients supplied by \eqref{eq:Qzero-jets}.

Equations \eqref{eq:modular-master}--\eqref{eq:critical-mod-jets},
together with the finite subtractions in \eqref{eq:root-bracket},
answer the modular-transport question at \(a=1\) in the indicated
order neighborhoods. General horizontal shifts, and finite
polylogarithmic reductions of the correction integrals, remain
separate questions.

